\documentclass[12pt]{article}
\usepackage[utf8]{inputenc}
\usepackage[T1]{fontenc}
\usepackage{hyperref}
\title{\'Ew\'e Language Resources for K-12 Curriculum Development}
\author{Compiled by Claude Sonnet 5.5 (claude-sonnet-5-5) for Andr\'as Kornai}
\date{}
\begin{document}
\maketitle

\section{Introduction}
\'Ew\'e (ISO 639-3 \texttt{ewe}) is spoken in Ghana, Togo and Benin. It is officially recognised in Ghana and taught in primary and secondary schools there, and in some secondary schools in Togo, where integration into education is reported as incomplete. This report is based on one round of searches and is not exhaustive.

\section{Primary Educational Resources}
\begin{itemize}
\item Omniglot overview of Ewe, including its writing and use in schools: \url{https://www.omniglot.com:443/writing/ewe.htm}
\item Indiana University CeLTOP Ewe language page: \url{https://celtop.indiana.edu/portal/languages/ewe/index.html}
\end{itemize}
Basic Ewe for foreign students (Dzablu-Kumah, 2015) and the Ewe Basic Course (1968) were mentioned in results, but no direct URL for them was seen.

\section{Linguistic and Reference Materials}
\begin{itemize}
\item Lexilogos Ewe dictionary links: \url{https://www.lexilogos.com/english/ewe_dictionary.htm}
\item Lexilogos French-language dictionary page: \url{https://www.lexilogos.com/ewe_dictionnaire.htm}
\item Global Voices interview on the 2024 French--Ewe dictionary \emph{Ele nam be mado Eʋəgbə}: \url{https://lingua.globalvoices.org/2025/01/interview-with-the-author-of-a-new-french-ewe-dictionary-in-togo/}
\end{itemize}

\section{Religious and Community Resources}
Digitised copy of the Ewe New Testament (SOAS): \url{https://digital.soas.ac.uk/AA00001744/00001/marc}

\section{Digital and Audio Resources}
No dedicated audio or software resources were confirmed in these searches.

\section{Curriculum Development Context}
Ewe has established school use in Ghana, so existing Ghanaian curriculum materials should be sought through Ghanaian education authorities (not located here). In Togo the language is a constitutionally recognised national language, but school integration is incomplete.
\end{document}